\scanpage{105}
\mspage{\unsure{204|104}}
Supposons $X\unsure{\gtrless|\bowtie}Y$, et soient $X'\ll X$, $Y'\ll Y$ avec
\[
X'_L=Y'_L=\varnothing.
\]
prouvons que $X'\parallel Y'$, i.e. \inserted{(prop.\,1)} que \deleted{pour} $\forall X''\leq X'$, $Y''\leq Y'$, on a $X''\mathrel{|\circ|}Y''$. Comme $\deleted{X''}X''_L\leq X'_L=\varnothing$ et \deleted{idem pour} \inserted{\unsure{de|idem} $\unsure{=|\equiv}$} $Y''_L=\varnothing$, OPS $X''=X'$, $Y''=Y'$. Soient $X_1\inserted{,Y_1}$ tels que \deleted{$X'\overset{\circ}{\ll}X_1\leq X,$}
\margin{$\unsure{Y}$}
\[
X'\overset{\circ}{\ll}X_1\leq X,\quad Y'\overset{\circ}{\ll}Y_1\leq Y
\]
On a $X_1\notin L$. (sinon on aurait $X'\in\unsure{\mathring L|L}$ et $X'_L\ne\varnothing$) et de $=$ $Y_1\notin L$, donc $X_1\mathrel{|\circ|}Y_1$ par hypo $X\gtrless Y$. Donc par Mag $5''$, on conclut $X'\mathrel{|\circ|}Y'$, qed

Supposons d'autre part l'hypothèse dite sur $X,Y$, en termes de $X',Y'$, satisfaite, prouvons que $X\gtrless Y$. Soit donc $X'\in\widetilde X\setminus L\deleted{\unsure{X}}$, $Y'=\widetilde Y\setminus L$, prouvons $X'\mathrel{|\circ|}Y'$.
% margin note written sideways / slanted; short note above the longer marginal paragraph
\margin{\unsure{condition} \unsure{on a ci-dessus}}
\deleted{\unsure{Ici} \unsure{on} \unsure{tout} \unsure{a priori} \unsure{soit}}
Sûrement, on a

% Two vertical emphasis strokes beside the lemma.
\margin{$\Vert$}
\uline{Lemme} La condition C sur $(X,Y)$ implique que pour $X',Y'$ comme dans l'énoncé de cette condition, et pour $x\in\operatorname{omb}(X')^{\circ}$, $y\in\operatorname{omb}(Y')^{\circ}$, on a $x\parallel y$.

% margin note written sideways / slanted
\margin{Si on suppose At L2$^+$, la condition C \inserted{pour $X,Y$} \unsure{n'est} ni plus, ni moins que : \deleted{$\operatorname{omb}(X)$} $\forall x\in\unsure{\operatorname{omb}(X)|\operatorname{omb}(X)^{\circ}}$, $y\in\unsure{\operatorname{omb}(Y)|\operatorname{omb}(Y)^{\circ}}$ \deleted{\unsure{et} $\unsure{\operatorname{omb}(y)}$}, on a $x\parallel y$ -- i.e. le cas particulier de C correspondant au cas $X'\in\unsure{\mathcal L|L}$, $Y'\in\mathcal L$.}

En effet, on aura bien sûr $\deleted{\unsure{x_L}}x_L=\varnothing\deleted{y_L=\varnothing}$ (sinon on aurait $x_L=x$ donc $x\ll L$ \deleted{donc} donc $\operatorname{omb}(X')\cap\operatorname{omb}(L)=\operatorname{omb}(X'_L)\ne\varnothing$, contrairement à $X'_L=\varnothing$), et de $=$ $y_L=\varnothing$, donc la condition.

Donc moralement, la validité de At C (\unsure{spac|spec}) signifie \deleted{la validité de} \inserted{celle de la} \deleted{que si $X,Y\in\unsure{\mathcal M|M}$, \illegible{1}} \deleted{L'axiome de} \inserted{on} \deleted{on ait} ait

\uline{Proposition} (?) Soient $X,Y\in\mathcal M$. \inserted{(ex-$X',Y'$)} Alors
\[
X\mathrel{|\circ|}Y\quad\Longleftrightarrow\quad\forall x\in\operatorname{omb}(X)^{\circ},\ y\in\operatorname{omb}(Y)^{\circ},\ \text{on a }x\parallel y.
\]
